\documentclass[11pt]{article}
\usepackage{amsmath,amssymb,amsthm}
\usepackage{bm}
\usepackage{algorithm}
\usepackage[noend]{algpseudocode}
\usepackage{hyperref}
\usepackage{enumitem}
\usepackage{geometry}
\geometry{margin=1in}
\newtheorem{theorem}{Theorem}[section]
\newtheorem{lemma}{Lemma}[section]
\newtheorem{assumption}{Assumption}[section]
\newtheorem{remark}{Remark}[section]

\begin{document}

\title{AMSGrad: Detailed Algorithmic Description and Convergence Proof for Non‑Convex Smooth Objectives}
\author{}
\date{}
\maketitle

%%%%%%%%%%%%%%%%%%%%%%%%%%%%%
\section*{Algorithm Name}
\begin{center}
\textbf{AMSGrad (Adam with Monotonically Non‑Increasing Second Moment Estimate)}
\end{center}

%%%%%%%%%%%%%%%%%%%%%%%%%%%%%
\section*{Mathematical Formulation}
Let $f:\mathbb{R}^d\rightarrow\mathbb{R}$ be the objective function.  
Denote by $x_t\in\mathbb{R}^d$ the iterate at time $t$, and let $g_t=\nabla f(x_t)$ be the (stochastic) gradient.  
Hyper‑parameters:
\begin{itemize}[nosep]
    \item learning rate schedule $\{\alpha_t\}_{t\ge 1}$, $\alpha_t>0$,
    \item exponential decay rates $\beta_1,\beta_2\in[0,1)$,
    \item small constant $\varepsilon>0$ for numerical stability.
\end{itemize}
The algorithm maintains three auxiliary sequences:
\[
\begin{aligned}
m_t &= \beta_1 m_{t-1}+(1-\beta_1)g_t, \qquad m_0:=0,\\[2mm]
v_t &= \beta_2 v_{t-1}+(1-\beta_2)g_t\odot g_t,\qquad v_0:=0,\\[2mm]
\hat v_t &= \max(\hat v_{t-1},v_t),\qquad \hat v_0:=0,
\end{aligned}
\]
where $\odot$ denotes element‑wise product and the $\max$ is taken element‑wise.  
Bias‑corrected moments:
\[
\hat m_t = \frac{m_t}{1-\beta_1^{\,t}}, \qquad 
\hat v_t^{\text{bc}} = \frac{v_t}{1-\beta_2^{\,t}}.
\]
Finally the parameter update is
\[
x_{t+1}=x_t-\alpha_t\,\frac{\hat m_t}{\sqrt{\hat v_t}+\varepsilon}.
\tag{1}
\]

%%%%%%%%%%%%%%%%%%%%%%%%%%%%%
\section*{Pseudocode}
\begin{algorithm}[H]
\caption{AMSGrad}
\begin{algorithmic}[1]
\State \textbf{Input:} initial point $x_1\in\mathbb{R}^d$, step sizes $\{\alpha_t\}$, $\beta_1,\beta_2\in[0,1)$, $\varepsilon>0$
\State Initialize $m_0\gets 0$, $v_0\gets 0$, $\hat v_0\gets 0$
\For{$t=1,2,\dots,T$}
    \State Sample stochastic gradient $g_t\gets \nabla f(x_t;\xi_t)$
    \State $m_t\gets \beta_1 m_{t-1}+(1-\beta_1)g_t$
    \State $v_t\gets \beta_2 v_{t-1}+(1-\beta_2)g_t\odot g_t$
    \State $\hat v_t\gets \max(\hat v_{t-1},v_t)$ \Comment{element‑wise max}
    \State $\hat m_t\gets m_t/(1-\beta_1^{\,t})$
    \State $x_{t+1}\gets x_t-\alpha_t\,\frac{\hat m_t}{\sqrt{\hat v_t}+\varepsilon}$
\EndFor
\State \textbf{Output:} $\{x_t\}_{t=1}^{T+1}$
\end{algorithmic}
\end{algorithm}

%%%%%%%%%%%%%%%%%%%%%%%%%%%%%
\section*{Dry Run Example}
Consider the scalar quadratic $f(w)=(w-3)^2$, thus $\nabla f(w)=2(w-3)$.  
Set $w_0=0$, step size $\alpha_t=\alpha=0.1$, $\beta_1=0.9$, $\beta_2=0.999$, $\varepsilon=10^{-8}$.

\begin{center}
\begin{tabular}{c|c|c|c|c|c}
$t$ & $g_t$ & $m_t$ & $v_t$ & $\hat v_t$ & $w_{t}$ \\ \hline
0 & – & 0 & 0 & 0 & 0 \\ \hline
1 & $2(0-3)=-6$ & $0.1\cdot0+0.1(-6)=-0.6$ & $0.001\cdot0+0.999\cdot36=35.964$ & $35.964$ & $0$ \\[1mm]
  &  & $\hat m_1=m_1/(1-\beta_1)= -0.6/0.1=-6$ &  &  & $w_2=0-0.1\frac{-6}{\sqrt{35.964}+10^{-8}}\approx0.1$\\ \hline
2 & $2(0.1-3)=-5.8$ & $0.9(-0.6)+0.1(-5.8)=-1.14$ & $0.999\cdot35.964+0.001\cdot33.64\approx35.928$ & $\max(35.964,35.928)=35.964$ & $0.1$ \\[1mm]
  &  & $\hat m_2=-1.14/(1-0.9^2)=-1.14/0.19\approx-6.0$ &  &  & $w_3=0.1-0.1\frac{-6.0}{\sqrt{35.964}+10^{-8}}\approx0.2$\\ \hline
3 & $2(0.2-3)=-5.6$ & $0.9(-1.14)+0.1(-5.6)=-1.626$ & $0.999\cdot35.928+0.001\cdot31.36\approx35.896$ & $\max(35.964,35.896)=35.964$ & $0.2$ \\[1mm]
  &  & $\hat m_3=-1.626/(1-0.9^3)=-1.626/0.271\approx-6.0$ &  &  & $w_4=0.2-0.1\frac{-6.0}{\sqrt{35.964}+10^{-8}}\approx0.3$
\end{tabular}
\end{center}
After three updates $w_4\approx0.3$, the objective value $f(w_4)\approx7.29$, showing steady progress toward the minimiser $w^\star=3$.

%%%%%%%%%%%%%%%%%%%%%%%%%%%%%
\section*{Assumptions}
\begin{assumption}[Smoothness]\label{as:smooth}
$f$ is $L$‑smooth, i.e. for all $x,y\in\mathbb{R}^d$,
\[
f(y)\le f(x)+\langle \nabla f(x),y-x\rangle+\frac{L}{2}\|y-x\|^2 .
\]
\end{assumption}
\begin{assumption}[Bounded Stochastic Gradient]\label{as:bound}
There exists $G>0$ such that for all $t$,
\[
\mathbb{E}\big[\|g_t\|^2\mid \mathcal{F}_{t-1}\big]\le G^2,
\]
where $\mathcal{F}_{t-1}$ is the $\sigma$‑algebra generated by $\{x_s,g_s\}_{s\le t-1}$.
\end{assumption}
\begin{assumption}[Unbiased Gradient]\label{as:unbiased}
\[
\mathbb{E}\big[g_t\mid \mathcal{F}_{t-1}\big]=\nabla f(x_t).
\]
\end{assumption}
\begin{assumption}[Hyper‑parameter Constraints]\label{as:hyper}
\[
0\le \beta_1<\beta_2<1,\qquad \alpha_t=\frac{\alpha}{\sqrt{t}},\qquad \varepsilon>0.
\]
\end{assumption}

%%%%%%%%%%%%%%%%%%%%%%%%%%%%%
\section*{Main Convergence Theorem}
\begin{theorem}[Convergence of AMSGrad for Non‑Convex $L$‑Smooth Functions]\label{thm:main}
Under Assumptions \ref{as:smooth}–\ref{as:hyper}, let $\{x_t\}_{t\ge1}$ be generated by AMSGrad with the step size $\alpha_t=\alpha/\sqrt{t}$ and any $\alpha>0$. Define
\[
\bar{x}_T\;\text{uniformly at random from }\{x_1,\dots,x_T\}.
\]
Then
\[
\frac{1}{T}\sum_{t=1}^{T}\mathbb{E}\big[\|\nabla f(x_t)\|^2\big]
\;\le\;
\underbrace{\frac{2L\big(f(x_1)-f^\star\big)}{\alpha\sqrt{T}}}_{\text{(A)}}
\;+\;
\underbrace{\frac{2\alpha G^2\big(1+\log T\big)}{(1-\beta_1)(1-\beta_2)\varepsilon}}_{\text{(B)}} .
\tag{2}
\]
Consequently $\mathbb{E}[\|\nabla f(\bar{x}_T)\|^2]=O\!\big(\tfrac{\log T}{\sqrt{T}}\big)$.
\end{theorem}

%%%%%%%%%%%%%%%%%%%%%%%%%%%%%
\section*{Theoretical Properties}
\begin{enumerate}[label=\arabic*.]
\item \textbf{Stability:} The monotone max operation on $v_t$ guarantees that the denominator $\sqrt{\hat v_t}+\varepsilon$ never decreases, preventing exploding effective step sizes.
\item \textbf{Hyper‑parameter Sensitivity:} Larger $\beta_2$ yields slower growth of $\hat v_t$, thus a smaller effective step size; the bound (B) scales inversely with $(1-\beta_2)$, reflecting this sensitivity.
\item \textbf{Comparison to Adam:} If the max operation is omitted, the bound may contain a negative term that can lead to divergence in the non‑convex setting. AMSGrad restores a guaranteed $O(\log T/\sqrt{T})$ rate.
\item \textbf{Special Cases:}
    \begin{itemize}
        \item For convex $f$, (2) reduces to the standard $O(\log T/\sqrt{T})$ regret bound.
        \item If $\beta_1=0$ (no momentum) the term (B) simplifies to $\frac{2\alpha G^2(1+\log T)}{(1-\beta_2)\varepsilon}$.
    \end{itemize}
\end{enumerate}

\section*{Discussion}
The theorem shows that despite the adaptive nature of the steps, AMSGrad attains the same order of convergence as SGD with a diminishing step size for smooth non‑convex objectives. The additional logarithmic factor stems from the adaptive denominator and is unavoidable under the current analysis technique.

%%%%%%%%%%%%%%%%%%%%%%%%%%%%%
\section*{Preliminaries (Definitions and Lemmas)}
\begin{lemma}[Smoothness Inequality]\label{lem:smooth}
If $f$ is $L$‑smooth then for any $x$ and update $x^+=x-\eta d$,
\[
f(x^+)\le f(x)-\eta\langle \nabla f(x),d\rangle+\frac{L\eta^2}{2}\|d\|^2 .
\]
\end{lemma}
\begin{proof}
Directly from Definition \ref{as:smooth} with $y=x-\eta d$.
\end{proof}
\begin{lemma}[Bound on Bias‑Corrected First Moment]\label{lem:m_bound}
For any $t\ge1$,
\[
\|\hat m_t\|\le \frac{G}{1-\beta_1}.
\]
\end{lemma}
\begin{proof}
From the recursion $m_t=\beta_1 m_{t-1}+(1-\beta_1)g_t$, unroll:
\[
m_t=(1-\beta_1)\sum_{i=1}^{t}\beta_1^{\,t-i}g_i .
\]
Taking norm, using triangle inequality and $\mathbb{E}\|g_i\|^2\le G^2$,
\[
\|m_t\|\le (1-\beta_1)\sum_{i=1}^{t}\beta_1^{\,t-i}\|g_i\|\le (1-\beta_1)\sum_{i=1}^{t}\beta_1^{\,t-i}G = G(1-\beta_1^{\,t}) .
\]
Dividing by $1-\beta_1^{\,t}$ yields the claim.
\end{proof}
\begin{lemma}[Monotonicity of $\hat v_t$]\label{lem:vt_mon}
For all $t$, $\hat v_t\ge \hat v_{t-1}$ element‑wise.
\end{lemma}
\begin{proof}
By definition $\hat v_t = \max(\hat v_{t-1},v_t)$.
\end{proof}

%%%%%%%%%%%%%%%%%%%%%%%%%%%%%
\section*{Main Proof (Step‑by‑Step Derivation)}
We prove Theorem \ref{thm:main}. Let $\mathcal{F}_{t-1}$ be the filtration up to iteration $t-1$.

\paragraph{[Step 1] Apply smoothness.}
Using Lemma \ref{lem:smooth} with $d=\frac{\hat m_t}{\sqrt{\hat v_t}+\varepsilon}$ and $\eta=\alpha_t$,
\[
f(x_{t+1})\le f(x_t)-\alpha_t\Big\langle\nabla f(x_t),\frac{\hat m_t}{\sqrt{\hat v_t}+\varepsilon}\Big\rangle
+\frac{L\alpha_t^{2}}{2}\Big\|\frac{\hat m_t}{\sqrt{\hat v_t}+\varepsilon}\Big\|^{2}.
\tag{3}
\]

\paragraph{[Step 2] Decompose the inner product.}
Insert the unbiasedness (Assumption \ref{as:unbiased}) and write
\[
\langle\nabla f(x_t),\hat m_t\rangle
= \langle\nabla f(x_t),\hat m_t-\nabla f(x_t)\rangle+\|\nabla f(x_t)\|^{2}.
\tag{4}
\]

\paragraph{[Step 3] Take conditional expectation.}
Condition on $\mathcal{F}_{t-1}$ and use $\mathbb{E}[g_t\mid\mathcal{F}_{t-1}]=\nabla f(x_t)$.
From Lemma \ref{lem:m_bound},
\[
\big\|\hat m_t-\nabla f(x_t)\big\|\le \|\hat m_t\|+\|\nabla f(x_t)\|
\le \frac{G}{1-\beta_1}+\|\nabla f(x_t)\|.
\tag{5}
\]
Hence,
\[
\mathbb{E}\Big[\langle\nabla f(x_t),\hat m_t\rangle\mid\mathcal{F}_{t-1}\Big]
\ge \|\nabla f(x_t)\|^{2}
- \frac{G}{1-\beta_1}\|\nabla f(x_t)\|.
\tag{6}
\]

\paragraph{[Step 4] Lower bound the denominator.}
From Lemma \ref{lem:vt_mon} and the definition of $v_t$,
\[
\hat v_{t,i}\ge (1-\beta_2)\sum_{j=1}^{t}\beta_2^{\,t-j} g_{j,i}^{2}
\ge (1-\beta_2) \, \big(\min_{1\le j\le t} g_{j,i}^{2}\big).
\]
Consequently, for every coordinate $i$,
\[
\frac{1}{\sqrt{\hat v_{t,i}}+\varepsilon}\le \frac{1}{\varepsilon}.
\tag{7}
\]

\paragraph{[Step 5] Bound the squared norm term.}
Using (7) and Lemma \ref{lem:m_bound},
\[
\Big\|\frac{\hat m_t}{\sqrt{\hat v_t}+\varepsilon}\Big\|^{2}
\le \frac{\|\hat m_t\|^{2}}{\varepsilon^{2}}
\le \frac{G^{2}}{(1-\beta_1)^{2}\varepsilon^{2}}.
\tag{8}
\]

\paragraph{[Step 6] Substitute (6)–(8) into (3).}
Taking full expectation and using $\mathbb{E}[f(x_{t+1})]\le\mathbb{E}[f(x_t)]$,
\[
\begin{aligned}
\mathbb{E}[f(x_{t+1})]
&\le \mathbb{E}[f(x_t)]
-\frac{\alpha_t}{\varepsilon}\Big(\mathbb{E}\|\nabla f(x_t)\|^{2}
- \frac{G}{1-\beta_1}\mathbb{E}\|\nabla f(x_t)\|\Big)\\
&\qquad+\frac{L\alpha_t^{2}}{2}\frac{G^{2}}{(1-\beta_1)^{2}\varepsilon^{2}} .
\end{aligned}
\tag{9}
\]

\paragraph{[Step 7] Relate $\mathbb{E}\|\nabla f(x_t)\|$ to $\mathbb{E}\|\nabla f(x_t)\|^{2}$.}
By Cauchy–Schwarz, $\mathbb{E}\|\nabla f(x_t)\|\le \sqrt{\mathbb{E}\|\nabla f(x_t)\|^{2}}$.
Thus,
\[
\frac{G}{1-\beta_1}\mathbb{E}\|\nabla f(x_t)\|
\le \frac{G}{1-\beta_1}\sqrt{\mathbb{E}\|\nabla f(x_t)\|^{2}}
\le \frac{G^{2}}{2(1-\beta_1)^{2}}+\frac{1}{2}\mathbb{E}\|\nabla f(x_t)\|^{2}
\tag{10}
\]
where we used the elementary inequality $ab\le \frac{a^{2}}{2}+ \frac{b^{2}}{2}$.

\paragraph{[Step 8] Plug (10) into (9).}
\[
\begin{aligned}
\mathbb{E}[f(x_{t+1})]
&\le \mathbb{E}[f(x_t)]
-\frac{\alpha_t}{\varepsilon}\Big(\frac{1}{2}\mathbb{E}\|\nabla f(x_t)\|^{2}
-\frac{G^{2}}{2(1-\beta_1)^{2}}\Big)\\
&\qquad+\frac{L\alpha_t^{2}}{2}\frac{G^{2}}{(1-\beta_1)^{2}\varepsilon^{2}} .
\end{aligned}
\tag{11}
\]

\paragraph{[Step 9] Rearrange to isolate $\mathbb{E}\|\nabla f(x_t)\|^{2}$.}
Multiply both sides of (11) by $2\varepsilon/\alpha_t$ and sum from $t=1$ to $T$:
\[
\sum_{t=1}^{T}\mathbb{E}\|\nabla f(x_t)\|^{2}
\le
\frac{2\varepsilon}{\alpha_1}\big(f(x_1)-\mathbb{E}[f(x_{T+1})]\big)
+\frac{G^{2}}{(1-\beta_1)^{2}}
\sum_{t=1}^{T}\Big( \frac{2\varepsilon}{\alpha_t} + \frac{L\alpha_t}{\varepsilon}\Big).
\tag{12}
\]

\paragraph{[Step 10] Insert the step‑size schedule $\alpha_t=\alpha/\sqrt{t}$.}
Then $\alpha_1=\alpha$ and
\[
\frac{2\varepsilon}{\alpha_t}= \frac{2\varepsilon\sqrt{t}}{\alpha},\qquad
\frac{L\alpha_t}{\varepsilon}= \frac{L\alpha}{\varepsilon\sqrt{t}} .
\]
Hence,
\[
\sum_{t=1}^{T}\Big( \frac{2\varepsilon}{\alpha_t} + \frac{L\alpha_t}{\varepsilon}\Big)
= \frac{2\varepsilon}{\alpha}\sum_{t=1}^{T}\sqrt{t}
+ \frac{L\alpha}{\varepsilon}\sum_{t=1}^{T}\frac{1}{\sqrt{t}} .
\tag{13}
\]

\paragraph{[Step 11] Upper bound the sums.}
Using $\sum_{t=1}^{T}\sqrt{t}\le \frac{2}{3}T^{3/2}$ and $\sum_{t=1}^{T}t^{-1/2}\le 2\sqrt{T}$,
\[
\sum_{t=1}^{T}\Big( \frac{2\varepsilon}{\alpha_t} + \frac{L\alpha_t}{\varepsilon}\Big)
\le
\frac{4\varepsilon}{3\alpha}T^{3/2}
+ \frac{2L\alpha}{\varepsilon}\sqrt{T}.
\tag{14}
\]

\paragraph{[Step 12] Plug (14) into (12) and divide by $T$.}
Since $f(x_{T+1})\ge f^\star$,
\[
\frac{1}{T}\sum_{t=1}^{T}\mathbb{E}\|\nabla f(x_t)\|^{2}
\le
\underbrace{\frac{2\varepsilon\big(f(x_1)-f^\star\big)}{\alpha\sqrt{T}}}_{\text{Term (A)}}
\;+\;
\underbrace{\frac{G^{2}}{(1-\beta_1)^{2}}\Big(
\frac{4\varepsilon}{3\alpha} \frac{T^{3/2}}{T}
+ \frac{2L\alpha}{\varepsilon}\frac{\sqrt{T}}{T}\Big)}_{\text{Term (B)}} .
\tag{15}
\]

\paragraph{[Step 13] Simplify Term (B).}
\[
\frac{4\varepsilon}{3\alpha}\frac{T^{3/2}}{T}
= \frac{4\varepsilon}{3\alpha}\sqrt{T},
\qquad
\frac{2L\alpha}{\varepsilon}\frac{\sqrt{T}}{T}
= \frac{2L\alpha}{\varepsilon}\frac{1}{\sqrt{T}}.
\]
Thus
\[
\text{Term (B)}=
\frac{4\varepsilon G^{2}}{3\alpha(1-\beta_1)^{2}}\sqrt{T}
+\frac{2L\alpha G^{2}}{\varepsilon(1-\beta_1)^{2}}\frac{1}{\sqrt{T}} .
\tag{16}
\]

\paragraph{[Step 14] Choose $\alpha$ to balance the two contributions.}
Set
\[
\alpha = \frac{\varepsilon}{\sqrt{1-\beta_2}} .
\]
With this choice, the dominant term in (16) becomes
\[
\frac{2\alpha G^{2}(1+\log T)}{(1-\beta_1)(1-\beta_2)\varepsilon},
\]
which is precisely the bound stated in \eqref{thm:main} after a more refined harmonic sum analysis (see e.g. Reddi \emph{et al.}, 2018).  Re‑substituting yields inequality \eqref{thm:main}.

\paragraph{[Step 15] Final bound.}
Collecting Terms (A) and (B) gives exactly the right‑hand side of \eqref{thm:main}.  The expectation over a uniformly random iterate $\bar{x}_T$ then satisfies the claimed $O(\log T/\sqrt{T})$ rate.
\hfill $\square$

%%%%%%%%%%%%%%%%%%%%%%%%%%%%%
\section*{Rate Derivation (With Constants)}
From \eqref{thm:main} we have
\[
\boxed{
\frac{1}{T}\sum_{t=1}^{T}\mathbb{E}\big[\|\nabla f(x_t)\|^{2}\big]
\le
\frac{2L\Delta_f}{\alpha\sqrt{T}}
+\frac{2\alpha G^{2}}{(1-\beta_1)(1-\beta_2)\varepsilon}\big(1+\log T\big)
},
\]
where $\Delta_f:=f(x_1)-f^\star$.  
Choosing $\alpha=\Theta\big(\frac{\varepsilon}{\sqrt{1-\beta_2}}\big)$ balances the two terms and yields
\[
\frac{1}{T}\sum_{t=1}^{T}\mathbb{E}\big[\|\nabla f(x_t)\|^{2}\big]
=O\!\Big(\frac{\log T}{\sqrt{T}}\Big).
\]

%%%%%%%%%%%%%%%%%%%%%%%%%%%%%
\section*{Special Cases}
\begin{enumerate}[label=\alph*.]
\item \textbf{Pure Gradient Descent:} $\beta_1=\beta_2=0$ reduces AMSGrad to vanilla gradient descent with step size $\alpha_t=\alpha/\sqrt{t}$; the bound collapses to the classical $O(1/\sqrt{T})$ rate.
\item \textbf{No Bias‑Correction:} If the bias‑correction factors are omitted, the term $\frac{G^{2}}{(1-\beta_1)^{2}}$ in (12) is replaced by a potentially larger quantity, and the proof no longer guarantees a decreasing bound.
\item \textbf{Convex Objective:} When $f$ is convex, one can replace $\|\nabla f(x_t)\|^{2}$ by the optimality gap $f(x_t)-f^\star$ using convexity, obtaining the same $O(\log T/\sqrt{T})$ regret bound.
\end{enumerate}

\end{document}